% 冥王星（综述）
% license CCBYSA3
% type Wiki

本文根据 CC-BY-SA 协议转载翻译自维基百科\href{https://en.wikipedia.org/wiki/Pluto}{相关文章}。

\begin{figure}[ht]
\centering
\includegraphics[width=6cm]{./figures/15a109b37bfc9d15.png}
\caption{} \label{fig_Pluto_1}
\end{figure}

















































\subsection{另见}
\begin{itemize}
\item 我如何“杀死”冥王星，以及它为什么罪有应得
\item 冥王星地质特征列表
\item 占星术中的冥王星
\item 小说中的冥王星
\item 太阳系行星的统计
\end{itemize}
\subsection{备注}\\
a.这张照片是由拉尔夫2015年7月14日，新视野号上的成像仪，距离35，445公里 (22，025英里)\\
b.这里的平均元素来自m é canique c é leste et de calcul des é ph é m é rides (IMCCE) 的外行星理论 (TOP2013) 解决方案。它们指的是标准的equinox J2000，太阳系的重心和纪元J2000。\\
c.由半径 \(r\) 计算表面积：\(4\pi r^{2}\)。\\
d.由半径 \(r\) 计算体积 \(v\)：\(4\pi r^{3}/3\)。\\
e.由质量 \(M\)、引力常数 \(G\) 和半径 \(r\) 计算表面重力：\(GM/r^{2}\)。\\
f.由质量 \(M\)、引力常数 \(G\) 和半径 \(r\) 计算逃逸速度：\(\sqrt{2GM/r}\)。\\
g.基于事实表中关于地球与冥王星最小与最大距离的几何关系以及冥王星半径\\
h.1919年一位法国天文学家曾为“X行星”建议名称“Pluto”，但没有迹象表明洛厄尔天文台的工作人员知道这一点\(^\text{[29]}\)。\\
i.例如，⟨♇⟩（Unicode: U+2647 ♇ PLUTO）出现在一篇发表于2004年的文章中，该文在2006年IAU定义发布之前，通过符号表识别行星\(^\text{[33]}\)；但在2016年的一张行星、矮行星与卫星示意图中，该符号并未出现，图中仅用符号表示IAU认定的八大行星\(^\text{[34]}\)。（一般来说，行星符号在天文学中并不常见，并且IAU不鼓励使用\(^\text{[35]}\)。）\\
j.bident符号 (U + 2BD3⯓冥王星形成两个) 自IAU决定矮行星以来也有一些天文用途，例如在NASA/JPL发布的关于矮行星的公共教育海报中黎明2015年的任务，其中IAU宣布的五个矮行星中的每一个都收到一个符号。\(^\text{[37]}\)此外，在占星术来源中还发现了冥王星的其他几个符号，\(^\text{[38]}\)包括Unicode接受的三个:,U + 2BD4⯔冥王星形成三,主要在南欧使用；/,U + 2BD6⯖冥王星形成五(以各种方向发现，显示冥王星的轨道横穿海王星的轨道)，主要在北欧使用; 和,U + 2BD5⯕冥王星形成四,用于天王星占星术。\(^\text{[39]}\) ⯔ ⯖⯖ ⯕\\
k.在其音韵学广泛不同于希腊语的,如索马里语 Buluuto和纳瓦霍人 Tłóotoo。\\
l.1978年Charon的发现使天文学家能够准确地计算出Plutonian系统的质量。但这并没有表明这两个物体的个体质量，只有在2005年底发现冥王星的其他卫星后才能估计。结果，由于冥王星于1989年来到近日点，因此大多数冥王星近日点日期估计都是基于冥王星-Charon重心。Charon来到近日点1989年9月4日。冥王星-Charon重心来到近日点1989年9月5日。冥王星来到近日点1989年9月8日。\\
m.由于冥王星轨道的偏心率，一些人认为它曾经是一个海王星的卫星。\(^\text{[102]}\)\\
n.矮行星厄里斯与冥王星的大小大致相同，约2330公里; 厄里斯的质量比冥王星大28\%。Eris是一个散乱圆盘对象,通常被认为是与冥王星等柯伊伯带天体截然不同的种群; 冥王星是柯伊伯带中最大的天体，不包括散落的圆盘天体。
\subsection{参考资料}
\begin{enumerate}
\item "Plutonian"。牛津英语词典(在线版)。牛津大学出版社。 (订阅或参与机构成员必需。)
\item 西蒙，J.L.；弗朗库，G.；Fienga, A.;Manche，H. (2013年9月)。“新的分析行星理论VSOP2013和TOP2013”(PDF)。天文学和天体物理学。557(2): a49。Bibcode:2013A & A...557A..49S。doi:10.1051/0004-6361/201321843。S2CID56344625。已存档(PDF)从原来的2022年5月27日。已检索2月26日，2024。   更清晰和常用格式的元素在电子表格中 已存档2016年5月15日，在回程机和原来的TOP2013元素在这里。 已存档2021年10月19日，在回程机
\item 威廉姆斯，大卫R.(2015年7月24日)。"冥王星概况介绍"。NASA.已存档从原来的2015年11月19日。已检索8月6日，2015。
\item “用于冥王星重心的地平线在线星历系统”。JPL Horizons在线星历系统@ 太阳能系统动力学组。已存档从2011年5月10日起。已检索1月16日，2011年。(观察者位置 @ 太阳，观察者位于太阳中心)
\item 弗朗西斯·尼莫等人 (2017年)。“来自新视野号图像的冥王星和卡戎的平均半径和形状”。伊卡洛斯。287:12-29。arXiv:1603.00821。Bibcode:2017Icar .. 287...12N。doi:10.1016/j.icarus.2016.06.027。S2CID44935431。 
\item 斯特恩，S.A.；格伦迪，W.；麦金农，W.B.；韦弗，H.A.；Young，洛杉矶.(2017)。“新视野之后的冥王星系统”。天文学和天体物理学年度评论。2018:357-392。arXiv:1712.05669。Bibcode:2018ARA & A .. 56 .. 357S。doi:10.1146/annurev-astro-081817-051935。ISSN0066-4146。S2CID119072504。  
\item 斯特恩，S.A.; 等人 (2015)。“冥王星系统: 新视野号探索的初步结果”。科学。350(6258) aad1815:249-352。arXiv:1510.07704。Bibcode:2015Sci。350.1815年。doi:10.1126/科学.aad1815。PMID26472913。S2CID1220226。  
\item 布罗佐维奇，玛丽娜; 雅各布森，罗伯特·A。(2024年5月8日)。“新地平线后冥王星卫星的轨道和质量”。天文杂志。167(256): 256。Bibcode:2024AJ....167..256B。doi:10.3847/1538-3881/ad39f0。
\item 塞利格曼，考特尼。“轮换周期和日长”。已存档从原来的2018年9月29日。已检索6月12日，2021。
\item “行星物理参数”。ssd.jpl.nasa.gov。存档自原来的2025年8月10日。已检索10月2日，2025。冥王星赤道半径1188.3(km) 恒星旋转周期-6.3872 (d)
\item Archinal, Brent A.;迈克尔·F·阿赫恩；爱德华·G·鲍威尔；阿尔伯特·R·康拉德；Consolmagno，Guy J.; 等人 (2010年)。“IAU制图坐标和旋转要素工作组的报告: 2009年”(PDF)。天体力学与动力天文学。109(2):101-135.Bibcode:2011 CeMDA.109..101A。doi:10.1007/s10569-010-9320-4。S2CID189842666。存档自原来的(PDF)2016年3月4日。已检索9月26日，2018。  
\item "AstDys (134340) 冥王星星历"。意大利比萨大学数学系。已存档从原来的2020年1月17日。已检索6月27日，2010。
\item "JPL小体数据库浏览器: 134340 Pluto"。已存档从原来的2017年2月18日。已检索9月29日，2022年。
\item 阿莫斯，乔纳森 (2015年7月23日)。新视野: 冥王星可能有 “氮冰川”"。BBC新闻。已存档从原来的2017年10月27日。已检索7月26日，2015。从阳光和无线电波通过Plutonian “空气” 的过程中可以看出，表面的压力仅为10微巴。
\item “冥王星大气中有一氧化碳”。Physorg.com。2011年4月19日已存档从2011年5月11日起。已检索11月22日，2011年。
\item 克莱德·汤博和帕特里克·摩尔 (2008)走出黑暗: 冥王星
\item 克罗斯韦尔，肯 (1997)。行星探索: 外星太阳系的史诗般的发现。纽约: 自由新闻。第43页。ISBN 978-0-684-83252-4。已存档从原来的2024年2月26日。已检索4月15日，2022年。
\item 汤博，克莱德·W。(1946)。寻找第九颗行星，冥王星太平洋天文学会传单。5(209):73-80。Bibcode:1946ASPL....5...73T。
\item Hoyt, William G. (1976). "W. H. Pickering's Planetary Predictions and the Discovery of Pluto". Isis. 67 (4): 551–564. doi:10.1086/351668. JSTOR 230561. PMID 794024. S2CID 26512655.
\item Littman, Mark (1990). Planets Beyond: Discovering the Outer Solar System. Wiley. p. 70. ISBN 978-0-471-51053-6.
\item Buchwald, Greg; Dimario, Michael; Wild, Walter (2000). "Pluto is Discovered Back in Time". Amateur–Professional Partnerships in Astronomy. 220. San Francisco: 335. Bibcode:2000ASPC..220..355B. ISBN 978-1-58381-052-1.
\item Croswell 1997, p. 50.
\item Croswell 1997, p. 52.
\item Rao, Joe (March 11, 2005). "Finding Pluto: Tough Task, Even 75 Years Later". Space.com. Archived from the original on August 23, 2010. Retrieved September 8, 2006.
\item Kevin Schindler & William Grundy (2018) Pluto and Lowell Observatory, pp. 73–79.
\item Croswell 1997, pp. 54–55.
\item Rincon, Paul (January 13, 2006). "The girl who named a planet". BBC News. Archived from the original on October 4, 2018. Retrieved April 12, 2007.
\item "Pluto Research at Lowell". Lowell Observatory. Archived from the original on April 18, 2016. Retrieved March 22, 2017.
\item Ferris (2012: 336) Seeing in the Dark
\item Scott & Powell (2018) The Universe as It Really Is
\item Coincidentally, as popular science author Martin Gardner and others have noted of the name "Pluto", "the last two letters are the first two letters of Tombaugh's name" Martin Gardner, Puzzling Questions about the Solar System (Dover Publications, 1997) p. 55
\item "NASA's Solar System Exploration: Multimedia: Gallery: Pluto's Symbol". NASA. Archived from the original on October 1, 2006. Retrieved November 29, 2011.
\item John Lewis，ed. (2004)。太阳系的物理和化学(2版)。爱思唯尔。第64页。
\item 陈晶晶; 大卫·基平 (2017)。“其他世界的质量和半径的概率预测”。天体物理学杂志。834(17)。美国天文学会: 8。arXiv:1603.08614。Bibcode:2017ApJ...834...17C。doi:10.3847/1538-4357/834/1/17。S2CID119114880。 
\item IAU风格手册(PDF)。1989。第27页。已存档(PDF)从2011年7月26日起。已检索1月29日，2022年。
\item Dane Rudhyar (1936)人格的占星术,将其归功于成立于1933年的Paul Clancy Publications。
\item NASA/JPL,什么是矮行星？ 已存档2021年12月8日，在回程机2015年4月22日
\item 弗雷德·盖廷斯 (1981)神秘，密封和炼金术符号词典。Routledge & Kegan Paul，伦敦。
\item 福克斯，大卫。"占星冥王星"(PDF)。www.unicode.org。Unicode。已存档(PDF)从原来的2020年11月12日。已检索10月1日，2021。
\item 海因里希斯，艾莉森·M。(2006)。"相形见绌"。匹兹堡论坛报-评论。存档自原来的2007年11月14日。已检索3月26日，2007。
\item 克拉克，大卫·L；大卫·霍巴特.(2000)。“关于传奇遗产的思考”(PDF)。已存档(PDF)从原来的2016年6月3日。已检索11月29日，2011年。
\item “行星语言学”。存档自原来的2007年12月17日。已检索6月12日，2007。
\item Renshaw，Steve; Ihara，Saori (2000)。“向侯ei Nojiri致敬”。存档自原来的2012年12月6日。已检索11月29日，2011年。
\item 浴袍.天王星，海王星和冥王星在中国，日本和越南。cjvlang.com。存档自原来的2011年7月20日。已检索11月29日，2011年。
\item 斯特恩，艾伦; 托伦，大卫·詹姆斯 (1997)。冥王星和Charon。亚利桑那大学出版社。pp. 206-208。ISBN 978-0-8165-1840-1。
\item Crommelin，Andrew Claude de la Cherois (1931)。“冥王星的发现”。皇家天文学会月刊。91(4):380-385。Bibcode:1931MNRAS..91..380。。doi:10.1093/mnras/91.4.380。
\item 尼科尔森，塞思·B.;Mayall，尼古拉斯·U。(1930年12月)。“冥王星质量的可能值”。太平洋天文学会的出版物。42(250): 350。Bibcode:1930PASP...42..350N。doi:10.1086/124071。
\item 尼科尔森，塞思·B.;Mayall，尼古拉斯·U。(1931年1月)。“冥王星的位置、轨道和质量”。天体物理学杂志。73: 1。Bibcode:1931ApJ。73。1N。doi:10.1086/143288。
\item 柯伊伯，杰拉德P.(1950)。“冥王星的直径”。太平洋天文学会的出版物。62(366):133-137。Bibcode:1950PASP...62..133K。doi:10.1086/126255。
\item Croswell 1997,第57页。
\item 詹姆斯·W·克里斯蒂；罗伯特·萨顿·哈灵顿(1978)。“冥王星的卫星”。天文杂志。83(8):1005-1008。Bibcode:1978AJ.....83.1005C。doi:10.1086/112284。S2CID120501620。 
\item 布伊，马克·W；威廉·M·格伦迪；Young，Eliot F.等人 (2006年)。“冥王星卫星的轨道和测光: Charon，S/2005 P1和S/2005 P2”。天文杂志。132(1):290-298。arXiv:astro-ph/0512491。Bibcode:2006AJ....132..290B。doi:10.1086/504422。S2CID119386667。 
\item Seidelmann，P。肯尼斯；罗伯特·萨顿·哈灵顿(1988)。“行星X-当前状态”。天体力学与动力天文学。43(1-4):55-68.Bibcode:1988年CeMec .. 43... 55s。doi:10.1007/BF01234554。S2CID189831334。 
\item 斯坦迪什，E.迈尔斯 (1993)。“行星X -- 光学观测中没有动力学证据”。天文杂志。105(5):200-2006。Bibcode:1993AJ....105.2000S。doi:10.1086/116575。
\item Standage，汤姆 (2000)。Neptune文件。企鹅.p. 168。ISBN 978-0-8027-1363-6。
\item 欧内斯特W.棕色，关于天王星扰动对海王星行星的预测 已存档2022年1月18日，在回程机,PNAS 1930年5月15日，16 (5) 364-371。
\item 尼尔·德格拉斯·泰森(2001年2月2日)。“天文学家回应冥王星非行星的说法”。Space.com。已存档从原来的2020年5月12日。已检索11月30日，2011年。
\item 梅茨格，菲利普·T.；塞克斯，马克·V。；斯特恩，艾伦; 鲁尼翁，柯比 (2019年)。“小行星从行星到非行星的重新分类”。伊卡洛斯。319:21-32。arXiv:1805.04115。Bibcode:2019Icar .. 319...21米。doi:10.1016/j.icarus.2018.08.026。S2CID119206487。 
\item 梅茨格，菲利普·T.；格伦迪，W。M、；马克·V·赛克斯；艾伦·斯特恩; 詹姆斯·F·贝尔三世；夏琳·德特里奇；鲁尼翁，柯比; 萨默斯，迈克尔 (2022)。“卫星是行星: 行星科学分类学中的科学有用性与文化目的论”。伊卡洛斯。374114768。arXiv:2110.15285。Bibcode:2022年Icar .. 37414768米。doi:10.1016/j.icarus.2021.114768。S2CID240071005。已存档从原来的2022年9月11日。已检索8月8日，2022年。 
\item “冥王星是一颗巨大的彗星吗？”。中央天文电报局。1999。已检索6月27日，2025。马斯登在1980年首次提出，“我们应该” [放弃] “第九行星” 的称谓，[分类] 冥王星 [...] 作为一个不寻常的小行星。
\item Moomaw，布鲁斯 (1999年6月28日)。“冥王星: 默认为行星”。spacedaily.com。已存档从2000年8月16日起。已检索6月27日，2025。
\item “冥王星的状态: 澄清”。国际天文学联合会,新闻稿。1999年2月3日。存档自原来的1999年4月22日
\item 布朗，迈克尔E。(2010)。我是如何杀死冥王星的，为什么它会来。ISBN 978-0-385-53108-5。
\item NASA资助的科学家发现了第十颗行星。NASA新闻稿。2005年7月29日已存档从原来的2020年5月12日。已检索2月22日，2007。
\item 索特，史蒂文 (2006年11月2日)。“什么是行星？”天文杂志。132(6):2513-2519。arXiv:astro-ph/0608359。Bibcode:2006AJ....132.2513S。doi:10.1086/508861。S2CID14676169。 
\item 国际宇航联合会。"IAF: Gonzalo Tancredi"。www.iafastro.org。已检索3月1日，2025。
\item "Cómo fue el día en que dos uruguayos lograron que Plutón dejara de ser considerado un planeta"。BBC新闻蒙多(西班牙语)。已检索3月1日，2025。
\item “天文学联盟2006年大会: 第5和第6号决议”(PDF)。IAU.2006年8月24日已存档(PDF)从原来的2009年6月20日。已检索6月15日，2008。
\item “IAU 2006年大会: IAU决议表决结果”。国际天文学联合会 (新闻稿-IAU0603)。2006年8月24日已存档从原来的2014年4月29日。已检索6月15日，2008。
\item Margot，jean-luc (2015)。“定义行星的定量标准”。天文杂志。150(6): 185。arXiv:1507.06300。Bibcode:2015AJ....150..185M。doi:10.1088/0004-6256/150/6/185。S2CID51684830。 
\item 索特，史蒂文 (2007)。“什么是行星？”。天文杂志。132(6)。美国自然历史博物馆天体物理系:2513-2519。arXiv:astro-ph/0608359。Bibcode:2006AJ....132.2513S。doi:10.1086/508861。S2CID14676169。已存档从原来的2013年11月6日。已检索4月9日，2012。 
\item 格林，丹尼尔·W.E.(2006年9月13日)。“(134340) 冥王星，(136199) Eris和 (136199) Eris I (Dysnomia)”(PDF)。IAU圆形(8747): 1。Bibcode:2006IAUC.8747....1G。存档自原来的2007年2月5日。已检索12月1日，2011年。
\item "JPL小体数据库浏览器"。加州理工学院。已存档从2011年7月21日起。已检索7月15日，2015。
\item 布里特，罗伯特·罗伊 (2006年8月24日)。冥王星降级: 不再是一颗极具争议的行星。Space.com。存档自原来的2010年12月27日。已检索9月8日，2006。
\item Ruibal，Sal (1999年1月6日)。天文学家质疑冥王星是否是真正的行星今日美国。
\item 布里特，罗伯特·罗伊 (2006年11月21日)。“为什么行星永远不会被定义”。Space.com。已存档从2009年5月24日的原始。已检索12月1日，2006。
\item 布里特，罗伯特·罗伊 (2006年8月24日)。科学家认为冥王星不再是一颗行星。NBC新闻。已存档从原来的2013年2月11日。已检索9月8日，2006。
\item 志贺，大卫 (2006年8月25日)。“新的行星定义引发了轩然大波”。NewScientist.com。已存档从原来的2010年10月3日。已检索9月8日，2006。
\item 克鲁格，杰弗里 (2025年2月18日)。“是什么让冥王星如此有趣”。时间。斯特恩认为矮行星的区别是荒谬的-对宇宙定义的任意解析。“小行星也是行星，” 他说。“仅仅因为太阳是一颗小恒星，我们就称它为矮星。我们不怕大量的行星; 我们不怕学童必须学习他们所有的名字。毕竟，孩子们不必记住周期表中的每个元素。
\item Buie, Marc W.(2006年9月)。“我对2006年IAU第5a和6a号决议的回应”。西南研究院。存档自原来的2007年6月3日。已检索12月1日，2011年。
\item 奥佛拜，丹尼斯 (2006年8月24日)。冥王星被降级为 “矮行星”"。纽约时报。已存档从原来的2022年6月22日。已检索12月1日，2011年。
\item 德沃尔，埃德娜 (2006年9月7日)。行星政治: 保护冥王星。Space.com。已存档从2011年8月4日起。已检索12月1日，2011年。
\item 霍顿，康斯坦斯 (2007年3月23日)。“修复冥王星”。科学。315(5819): 1643。doi:10.1126/science.315.5819.1643c。S2CID220102037。 
\item 古铁雷斯，乔尼·玛丽 (2007年)。“联合纪念碑。宣布冥王星为行星，并在立法机关宣布2007年3月13日为 “冥王星行星日”。。新墨西哥州立法机关。已存档从原来的2020年5月11日。已检索9月5日，2009年。
\item “伊利诺伊州大会: 第96届大会SR0046的法案状态”。ilga.gov。伊利诺伊州大会。已存档从2011年5月14日起。已检索3月16日，2011年。
\item “冥王星仍然是同一个冥王星”。独立报纸。美联社。2006年10月21日。已存档从原来的2017年12月1日。已检索11月29日，2011年。米老鼠有一只可爱的狗。
\item "“Plutoed” 被选为 “06年度最佳单词”。美联社。2007年1月8日已存档从2013年3月1日起。已检索1月10日，2007。
\item 桑切斯，卡梅伦 (2024年4月6日)。“冥王星再次成为一颗行星-至少在亚利桑那州”。npr.org。NPR。已检索4月12日，2024。
\item 明克尔，J.R。(2008年4月10日)。“重新点燃冥王星行星辩论是个好主意吗？”。科学美国人。已存档从2011年8月11日起。已检索12月1日，2011年。
\item 伟大的行星辩论: 科学作为过程。科学会议和教育家研讨会”。gpd.jhuapl.edu。约翰霍普金斯大学应用物理实验室。2008年6月27日。已存档从2011年8月17日起。已检索12月1日，2011年。
\item “科学家辩论行星定义并同意不同意”，行星科学研究所2008年9月19日的新闻稿，PSI.edu 已存档2011年7月15日，在回程机
\item “Plutoid被选为像冥王星这样的太阳系物体的名称”。巴黎:国际天文学联合会(新闻稿-IAU0804)。2008年6月11日。已存档从2011年8月10日起。已检索12月1日，2011年。
\item “Plutoids加入了太阳能家族”，《发现》杂志，2009年1月，第76页
\item 科学新闻，2008年7月5日，第7页
\item “冥王星将成为最遥远的星球”。JPL/NASA。1999年1月28日存档自原来的2010年9月2日。已检索1月16日，2011年。
\item 杰拉尔德·杰伊·萨斯曼; 杰克·智慧 (1988)。“冥王星运动混乱的数字证据”。科学。241(4864):433-437。Bibcode:1988Sci...241 .. 433S。doi:10.1126/science.241.4864.433。高密度脂蛋白:1721.1/6038。PMID17792606。S2CID1398095。存档自原来的2017年9月24日。已检索5月16日，2018。  
\item 智慧，杰克; 霍尔曼，马修 (1991)。“n体问题的辛映射”。天文杂志。102:1528-1538。Bibcode:1991AJ....102.1528W。doi:10.1086/115978。已存档从原始的2021年7月10日。已检索10月18日，2021。
\item 詹姆斯·G·威廉姆斯；本森，G。S。(1971)。“海王星-冥王星系统中的共振”。天文杂志。76: 167。Bibcode:1971AJ.....76..167W。doi:10.1086/111100。S2CID120122522。 
\item 万晓生; 黄天义; Innanen，Kim A。(2001)。“冥王星运动中的1:1超共振”。天文杂志。121(2):1155-1162。Bibcode:2001AJ....121.1155W。doi:10.1086/318733。
\item 亨特，麦克斯韦W.(2004)。“整个太阳系的无人科学探索”。空间科学评论。6(5): 501。Bibcode:1967年SSRv .... 6 .. 601小时。doi:10.1007/BF00168793。S2CID125982610。 
\item Malhotra，Renu (1997年)。“冥王星的轨道”。已存档从原始的2019年7月31日。已检索3月26日，2007。
\item 卡尔·萨根&安·德鲁扬(1997)。彗星。纽约: 兰登书屋。第223页。ISBN 978-0-3078-0105-0。已存档从原来的2024年2月26日。已检索10月18日，2021。
\item 的冥王星的黄道经度 已存档2024年2月13日，在回程机和海王星的 已存档2024年2月13日，在回程机可从JPL Horizons在线星历系统。
\item 阿尔夫文，汉尼斯; 阿伦尼乌斯，古斯塔夫 (1976)。“太阳系的SP-345演化”。已存档从原来的2007年5月13日。已检索3月28日，2007。
\item 科恩，C。J.;哈伯德，E。C.(1965)。“冥王星接近海王星的解放”。天文杂志。70: 10。Bibcode:1965AJ.....70...10C。doi:10.1086/109674。
\item 福雷，冈特; 梅辛，特蕾莎·M。(2007)。“冥王星和夏隆: 奇怪的一对”。行星科学导论。斯普林格.pp. 401-408。doi:10.1007/978-1-4020-5544-7。ISBN 978-1-4020-5544-7。
\item 吉姆·舒姆伯特；俄勒冈大学天文学121讲义 已存档2011年7月23日，在回程机,冥王星方向图 已存档2009年3月25日，在回程机
\item 约瑟夫·L·柯希维克；罗伯特·L·里珀丹；埃文斯，大卫A.(1997年7月25日)。“通过惯性交换真正的极地漂移对寒武纪早期大陆团块进行大规模重组的证据”。科学。277(5325):541-545。doi:10.1126/science.277.5325.541。ISSN0036-8075。S2CID177135895。  
\item 基恩，詹姆斯·T；松山，Isamu; Kamata，Shunichi; Steckloff，Jordan K.(2016)。“由于人造卫星Planitia内的挥发性负荷，冥王星的重新定向和断层”。自然。540(7631):90-93。Bibcode:2016 Natur.540. .. 90K。doi:10.1038/nature20120。PMID27851731。S2CID4468636。  
\item 托比亚斯·欧文；特德·L·鲁什；Cruikshank，Dale P.; 等人 (1993年)。“冥王星的表面冰和大气成分”。科学。261(5122):745-748。Bibcode:1993Sci。261。。745O。doi:10.1126/science.261.5122.745。JSTOR2882241。PMID17757212。S2CID6039266。   
\item 格伦迪，W.M.；奥尔金，c.b.；年轻，洛杉矶；Buie, M.W.;Young, E.F.(2013)。“冥王星冰的近红外光谱监测: 空间分布和长期演变”(PDF)。伊卡洛斯。223(2):710-721。arXiv:1301.6284。Bibcode:2013年Icar。。223。710克。doi:10.1016/j.icarus.2013.01.019。S2CID26293543。存档自原来的(PDF)2015年11月8日  
\item 纳迪亚·德雷克(2015年11月9日)。“冥王星上的浮山-你不能弥补这些东西”。国家地理。存档自原来的2015年11月13日。已检索12月23日，2016。
\item 布伊，马克·W；威廉·M·格伦迪；杨，艾略特·F; 等人 (2010年)。“冥王星和Charon与哈勃太空望远镜: I.从光曲线监测全球变化和改进的表面特性”。天文杂志。139(3):1117-1127。Bibcode:2010AJ....139.1117B。CiteSeerX10.1.1.625.7795。doi:10.1088/0004-6256/139/3/1117。S2CID1725219。已存档从原来的2015年7月20日。已检索2月10日，2010。 
\item Buie, Marc W."冥王星地图信息"。存档自原来的2011年6月29日。已检索2月10日，2010。
\item 雷·维拉德; 布伊，马克·W。(2010年2月4日)。“冥王星的新哈勃地图显示了表面变化”。新闻发布编号: STScI-2010-06。已存档从2016年9月1日起。已检索2月10日，2010。
\item 布伊，马克·W；威廉·M·格伦迪；杨，艾略特·F; 等人 (2010年)。“冥王星和Charon与哈勃太空望远镜: II。解决冥王星表面的变化和Charon的地图“。天文杂志。139(3):1128-1143。Bibcode:2010AJ....139.1128B。CiteSeerX10.1.1.625.7795。doi:10.1088/0004-6256/139/3/1128。S2CID9343680。已存档从原来的2015年7月7日。已检索2月10日，2010。  
\item 艾米丽·拉克达瓦拉(2016年10月26日)。“DPS/EPSC在冥王星系统及以后的新视野上的更新”。行星社会。已存档从原来的2018年10月8日。已检索10月26日，2016。
\item 麦金农，W.B.;尼莫，F.；黄，T.；Schenk, P.M、；白色，O.等。(2016年6月1日)。“富含挥发性氮冰的层中的对流驱动了冥王星的地质活力”。自然。534(7605):82-85。arXiv:1903.05571。Bibcode:2016 Natur.534. .. 82M。doi:10.1038/自然18289。PMID27251279。S2CID30903520。  
\item 特罗布里奇，A。J.;Melosh, H. J.;斯特克洛夫，J。K.;Freed, A.M。(2016年6月1日)。“剧烈对流作为冥王星多边形地形的解释”。自然。534(7605):79-81。Bibcode:2016 Natur.534. .. 79T。doi:10.1038/自然18016。PMID27251278。S2CID6743360。  
\item 艾米丽·拉克达瓦拉(2015年12月21日)。“来自AGU和DPS的冥王星更新: 来自混乱世界的漂亮图片”。行星社会。已存档从原来的2015年12月24日。已检索1月24日，2016。
\item Umurhan, O.(2016年1月8日)。在冥王星冰冻的 “心脏” 上探测神秘的冰川流"。blogs.nasa.gov。NASA.已存档从原来的2016年4月19日。已检索1月24日，2016。
\item Marchis, F.;特里林，D。E.(2016年1月20日)。“冥王星人造卫星的表面年龄必须小于1000万年”。PLOS ONE。11(1) e0147386。arXiv:1601.02833。Bibcode:2016PLoSO .. 1147386T。doi:10.1371/期刊.pone.0147386。PMC4720356。PMID26790001。  
\item 布勒，P。B.;英格索尔，A.P。(2017年3月23日)。“升华坑分布表明，冥王星Sputnik Planitia的对流单元表面速度为每年约10厘米”(PDF)。第48届月球与行星科学会议。已存档(PDF)从原来的2017年8月13日。已检索5月11日，2017。
\item Telfer, Matt W; Parteli, Eric J. R; Radebaugh, Jani; Beyer, Ross A; Bertrand, Tanguy; Forget, François; Nimmo, Francis; Grundy, Will M; Moore, Jeffrey M; Stern, S. Alan; Spencer, John; Lauer, Tod R; Earle, Alissa M; Binzel, Richard P; Weaver, Hal A; Olkin, Cathy B; Young, Leslie A; Ennico, Kimberly; Runyon, Kirby (2018). "Dunes on Pluto" (PDF). Science. 360 (6392): 992–997. Bibcode:2018Sci...360..992T. doi:10.1126/science.aao2975. PMID 29853681. S2CID 44159592. Archived (PDF) from the original on October 23, 2020. Retrieved April 12, 2020.
\item Robbins, Stuart J.; Dones, Luke (December 2023). "Impact Crater Databases for Pluto and Charon, Version 2". The Planetary Science Journal. 4 (12): 6. Bibcode:2023PSJ.....4..233R. doi:10.3847/PSJ/acf7be. S2CID 266147862. 233.
\item Hussmann, Hauke; Sohl, Frank; Spohn, Tilman (November 2006). "Subsurface oceans and deep interiors of medium-sized outer planet satellites and large trans-neptunian objects". Icarus. 185 (1): 258–273. Bibcode:2006Icar..185..258H. doi:10.1016/j.icarus.2006.06.005. Archived (PDF) from the original on August 31, 2015. Retrieved October 25, 2018.
\item "The Inside Story". pluto.jhuapl.edu – NASA New Horizons mission site. Johns Hopkins University Applied Physics Laboratory. 2007. Archived from the original on May 16, 2008. Retrieved February 15, 2014.
\item Overlooked Ocean Worlds Fill the Outer Solar System Archived December 26, 2018, at the Wayback Machine. John Wenz, Scientific American. October 4, 2017.
\item Samantha Cole. "An Incredibly Deep Ocean Could Be Hiding Beneath Pluto's Icy Heart". Popular Science. Archived from the original on September 27, 2016. Retrieved September 24, 2016.
\item 拉比，帕森特 (2020年6月22日)。“新的证据表明冥王星有些奇怪和令人惊讶的东西-这些发现将使科学家重新思考柯伊伯带天体的可居住性”。逆。已存档从原来的2020年6月23日。已检索6月23日，2020。
\item 比尔森，卡弗; 等人 (2020年6月22日)。“冥王星热启动和早期海洋形成的证据”。自然地球科学。769(7):468-472。Bibcode:2020年自然 .. 13 .. 468B。doi:10.1038/s41561-020-0595-0。S2CID219976751。已存档从原来的2020年6月22日。已检索6月23日，2020。 
\item 歌手，Kelsi N.(2022年3月29日)。“冥王星上的大规模低温火山表面”。自然通讯。13(1) 1542。arXiv:2207.06557。Bibcode:2022年纳科。。13.1542年代。doi:10.1038/s41467-022-29056-3。PMC8964750。PMID35351895。  
\item 约翰·戴维斯 (2001)。"超越冥王星 (摘录)"(PDF)。爱丁堡皇家天文台。存档自原来的(PDF)2011年7月15日。已检索3月26日，2007。
\item 冥王星有多大？“新视野解决了长达数十年的辩论”。NASA.2015年7月13日存档自原来的2017年7月1日。已检索7月13日，2015。
\item “冥王星和卡戎 | 天文学”。courses.lumenlearning.com。已存档从原来的2022年3月24日。已检索4月6日，2022年。长期以来，人们认为冥王星的质量与地球相似，因此被归类为第五类地行星，不知何故，在太阳系的最外围放错了地方。然而，还有其他异常现象，因为冥王星的轨道比任何其他行星都更加偏心，并且更倾向于我们太阳系的平面。只有在1978年发现冥王星的月球Charon之后，才可以测量冥王星的质量，结果发现它远小于地球的质量。
\item 关闭，莱尔德·M；威廉·J·梅林；大卫·J·托伦等人 (2000年)。“冥王星-卡隆的自适应光学成像和小行星45 Eugenia周围的月球的发现: 自适应光学在行星天文学中的潜力”。在Wizinowich，彼得L. (ed.)。自适应光学系统技术。第4007卷。pp。 787-795。Bibcode:2000SPIE.4007..787C。doi:10.1117/12.390379。S2CID122678656。 
\item 杨，艾略特·F；年轻，莱斯利A；Buie, Marc W.(2007)。“冥王星的半径”。美国天文学会，DPS第39号会议，#62.05; 美国天文学会公报。39: 541。Bibcode:2007DPS....39.6205Y。
\item 布朗，迈克尔E。(2010年11月22日)。“冥王星到底有多大？”。迈克·布朗的行星。已存档从2011年7月21日起。已检索6月9日，2015。 (2010年11月8日Franck Marchis)[死链接]
\item Lellouch，伊曼纽尔; de Bergh，凯瑟琳; Sicardy，布鲁诺; 等人 (2015年1月15日)。“探索冥王星大气中甲烷的空间、时间和垂直分布”。伊卡洛斯。246:268-278。arXiv:1403.3208。Bibcode:2015Icar .. 246 .. 268L。doi:10.1016/j.icarus.2014.03.027。S2CID119194193。 
\item Lakdawalla，Emily (2015年7月13日)。“冥王星减去一天: 冥王星遇到科学成果的第一个新视野”。行星社会。已存档从原来的2020年3月2日。已检索7月13日，2015。
\item NASA的新视野团队揭示了冥王星的新科学发现。NASA.2015年7月24日事件发生在52:30。存档自原来的2021年10月28日。已检索7月30日，2015。我们有一个范围可能超过70公里的不确定性，我们已经把它降到了正负2，它的中心是1186
\item “冥王星的条件: 令人难以置信的朦胧与流动的冰”。纽约时报。2015年7月24日已存档从原来的2015年7月28日。已检索7月24日，2015。
\item Croswell，Ken (1992)。“冥王星大气中的氮”。KenCroswell.com。新科学家。已存档从原来的2020年5月11日。已检索4月27日，2007。
\item 奥尔金，c.b.；年轻，洛杉矶；Borncamp，D.; 等人 (2015年1月)。“冥王星的大气层不会因包括2013年5月4日事件在内的掩星而崩溃的证据”。伊卡洛斯。246:220-225。Bibcode:2015年Icar .. 246 .. 220O。doi:10.1016/j.icarus.2014.03.026。高密度脂蛋白:10261/167246。已存档从原始的2021年9月29日。已检索9月8日，2017。
\item 凯利·比蒂 (2016)。“冥王星的大气层让研究人员感到困惑”。天空和望远镜。已存档从原来的2016年4月7日。已检索4月2日，2016。
\item 比，Ker (2006)。“天文学家: 冥王星比预期的要冷”。Space.com。已存档从原来的2012年10月19日。已检索11月30日，2011年-通过CNN.
\item 格莱斯顿，G。R、；斯特恩，S.A.;恩尼科，K.; 等人 (2016年3月)。“新视野号观测到的冥王星大气层”(PDF)。科学。351(6279) aad8866。arXiv:1604.05356。Bibcode:2016Sci...351.8866克。doi:10.1126/科学.aad8866。PMID26989258。S2CID32043359。存档自原来的(PDF)2016年5月21日。已检索6月12日，2016。    (补充材料)
\item “冥王星的大气层发生了什么？”。2020年5月22日。已存档从原始的2021年10月24日。已检索10月7日，2021。
\item "SwRI Scientists Confirm Decrease In Pluto's Atmospheric Density". Southwest Research Institute. October 4, 2021. Archived from the original on October 15, 2021. Retrieved October 7, 2021.
\item Lellouch, Emmanuel; Sicardy, Bruno; de Bergh, Catherine; et al. (2009). "Pluto's lower atmosphere structure and methane abundance from high-resolution spectroscopy and stellar occultations". Astronomy and Astrophysics. 495 (3): L17 – L21. arXiv:0901.4882. Bibcode:2009A&A...495L..17L. doi:10.1051/0004-6361/200911633. S2CID 17779043.
\item Gugliotta, Guy (November 1, 2005). "Possible New Moons for Pluto". The Washington Post. Archived from the original on October 20, 2012. Retrieved October 10, 2006.
\item "NASA's Hubble Discovers Another Moon Around Pluto". NASA. July 20, 2011. Archived from the original on May 12, 2020. Retrieved July 20, 2011.
\item Wall, Mike (July 11, 2012). "Pluto Has a Fifth Moon, Hubble Telescope Reveals". Space.com. Archived from the original on May 14, 2020. Retrieved July 11, 2012.
\item Buie, M.; Tholen, D.; Grundy, W. (2012). "The Orbit of Charon is Circular" (PDF). The Astronomical Journal. 144 (1): 15. Bibcode:2012AJ....144...15B. doi:10.1088/0004-6256/144/1/15. S2CID 15009477. Archived from the original (PDF) on April 12, 2020.
\item Showalter, M.R.; Hamilton, D.P. (June 3, 2015). "Resonant interactions and chaotic rotation of Pluto's small moons". Nature. 522 (7554): 45–49. Bibcode:2015Natur.522...45S. doi:10.1038/nature14469. PMID 26040889. S2CID 205243819.
\item 斯特恩，S.艾伦; 韦弗，哈罗德·A。Jr.; Steffl，Andrew J.; 等人 (2005)。“冥王星四重系统的特征和起源”。arXiv:astro-ph/0512599。
\item 亚历山德拉·维兹 (2015)。“冥王星的卫星同步移动”。自然。doi:10.1038/自然2015.17681。S2CID134519717。 
\item Matson, J.(2012年7月11日)。“冥王星的新月: 哈勃望远镜发现了第五颗普卢顿卫星”。科学美国人网站。已存档从原来的2016年8月31日。已检索7月12日，2012。
\item 德里克·理查森；凯文·J·沃尔什.(2005)。“二进制小行星”。地球与行星科学年度评论。34(1):47-81。Bibcode:2006AREPS..34...47R。doi:10.1146/annurev.earth.32.101802.120208。S2CID1692921。 
\item Sicardy，Bruno; Bellucci，aur é lie; Gendron，é ric; 等人 (2006年)。“查伦的大小和来自恒星掩星的大气层上限”。自然。439(7072):52-54.Bibcode:2006年自然439 .. 52s。doi:10.1038/nature04351。高密度脂蛋白:11336/39754。PMID16397493。S2CID4411478。  
\item 萨卡茨，R.；吻，Cs。；奥尔蒂斯，J.L.；莫拉莱斯，N.；Pál, A.;M ü ller，T.G.; 等人 (2023)。“从长期的地面和空间测光法中发现的矮行星 (136199) Eris的潮汐锁定旋转”。天文学与天体物理学。L3: 669。arXiv:2211.07987。Bibcode:2023A & A...669L...3S。doi:10.1051/0004-6361/202245234。S2CID253522934。 
\item 年轻的Leslie A.(1997)。《曾经和未来的冥王星》。科罗拉多州博尔德西南研究所。已存档从原来的2004年3月30日。已检索3月26日，2007。
\item NASA的哈勃发现冥王星的卫星在绝对混乱中翻滚。2015年6月3日。已存档从原来的2020年4月6日。已检索6月3日，2015。
\item de la Fuente Marcos，卡洛斯; de la Fuente Marcos，劳尔 (2012)。"Plutino 15810 (1994 JR1)，冥王星的一颗偶然的准卫星“。皇家天文学会的每月通知。427(1): l85。arXiv:1209.3116。Bibcode:2012MNRAS.427L..85D。doi:10.1111/j.1745-3933.2012.01350.x。S2CID118570875。 
\item 《冥王星的假月亮》。天空和望远镜。2012年9月24日已存档从原来的2020年2月20日。已检索9月24日，2012。
\item “新视野号收集了关于后冥王星天体的第一门科学”。NASA.2016年5月13日.存档自原来的2016年6月7日。已检索6月5日，2016。
\item 德拉富恩特·马科斯，卡洛斯; 德拉富恩特·马科斯，劳尔 (2016)。“分析标准: 偶然的准卫星确实是真正的准卫星”。皇家天文学会月刊。462(3):3344-3349。arXiv:1607.06686。Bibcode:2016MNRAS.462.3344D。doi:10.1093/mnras/stw1833。S2CID119284843。 
\item 波特，西蒙·B; 等人 (2016)。“KBO的第一个高相位观测: 来自柯伊伯带的 (15810) 1994 JR1的新视野成像”。天体物理学杂志快报。828(2): l15。arXiv:1605.05376。Bibcode:2016ApJ...828L..15P。doi:10.3847/2041-8205/828/2/L15。S2CID54507506。 
\item 杰拉德·柯伊伯 (1961)。行星和卫星。芝加哥: 芝加哥大学出版社。第576页。
\item 斯特恩，S.艾伦；大卫·J·托伦。(1997)。冥王星和Charon。亚利桑那大学出版社。第623页。ISBN 978-0-8165-1840-1. Archived from the original on February 26, 2024. Retrieved October 23, 2015.
\item Sheppard, Scott S.; Trujillo, Chadwick A.; Udalski, Andrzej; et al. (2011). "A Southern Sky and Galactic Plane Survey for Bright Kuiper Belt Objects". Astronomical Journal. 142 (4): 98. arXiv:1107.5309. Bibcode:2011AJ....142...98S. doi:10.1088/0004-6256/142/4/98. S2CID 53552519.
\item "Colossal Cousin to a Comet?". pluto.jhuapl.edu – NASA New Horizons mission site. Johns Hopkins University Applied Physics Laboratory. Archived from the original on November 13, 2014. Retrieved February 15, 2014.
\item Tyson, Neil deGrasse (1999). "Pluto Is Not a Planet". The Planetary Society. Archived from the original on September 27, 2011. Retrieved November 30, 2011.
\item Philip Metzger (April 13, 2015). "Nine Reasons Why Pluto Is a Planet". Philip Metzger. Archived from the original on April 15, 2015.
\item Wall, Mike (May 24, 2018). "Pluto May Have Formed from 1 Billion Comets". Space.com. Archived from the original on May 24, 2018. Retrieved May 24, 2018.
\item Glein, Christopher R.; Waite, J. Hunter Jr. (May 24, 2018). "Primordial N2 provides a cosmochemical explanation for the existence of Sputnik Planitia, Pluto". Icarus. 313 (2018): 79–92. arXiv:1805.09285. Bibcode:2018Icar..313...79G. doi:10.1016/j.icarus.2018.05.007. S2CID 102343522.
\item "Neptune's Moon Triton". The Planetary Society. Archived from the original on December 10, 2011. Retrieved November 30, 2011.
\item 戈麦斯R.美国；Gallardo T.;费尔南德斯J.A.;布鲁尼尼A.(2005)。“关于高近日点散射盘的起源: Kozai机制和平均运动共振的作用”。天体力学与动力天文学。91(1-2):109-129。Bibcode:2005年CeMDA。。91。109克。doi:10.1007/s10569-004-4623-y。高密度脂蛋白:11336/38379。S2CID18066500。  
\item 哈恩，约瑟夫·M。(2005)。“海王星迁移到搅动的柯伊伯带: 模拟与观测的详细比较”(PDF)。天文杂志。130(5):2392-2414。arXiv:astro-ph/0507319。Bibcode:2005AJ....130.2392H。doi:10.1086/452638。S2CID14153557。已存档(PDF)从2011年7月23日起。已检索3月5日，2008。  
\item 哈罗德·F·利维森；Morbidelli，Alessandro; Van Laerhoven，Christa; 等人 (2007年)。“天王星和海王星轨道动态不稳定期间柯伊伯带结构的起源”。伊卡洛斯。196(1):258-273.arXiv:0712.0553。Bibcode:2008Icar .. 196 .. 258L。doi:10.1016/j.icarus.2007.11.035。S2CID7035885。 
\item Malhotra，Renu (1995)。“冥王星轨道的起源: 对海王星以外太阳系的影响”。天文杂志。110: 420。arXiv:astro-ph/9504036。Bibcode:1995AJ....110..420M。doi:10.1086/117532。S2CID10622344。 
 “本月冥王星的视星等为m = 14.1。我们能看到它的焦距为3400毫米的11英寸反射镜吗？。新加坡科学中心。2002。存档自原来的2005年11月11日。已检索11月29日，2011年。
 “星期五如何在夜空中寻找冥王星”。Space.com。2014年7月3日已存档从原来的2022年4月6日。已检索4月6日，2022年。
 杨，艾略特·F；理查德·P·宾泽尔；克雷恩，基南 (2001)。“冥王星亚Charon半球的双色地图”。天文杂志。121(1):552-561。Bibcode:2001AJ....121..552Y。doi:10.1086/318008。
 布伊，马克·W；大卫·J·索伦；基思·霍恩 (1992)。“冥王星和Charon的反照率图: 初始相互事件结果”。伊卡洛斯。97(2):221-227。Bibcode:1992年伊卡尔。97。。211B。doi:10.1016/0019-1035(92)90129-U。已存档从2011年6月22日的原件。已检索2月10日，2010。
 Buie, Marc W.“冥王星地图是如何制作的”。存档自原来的2010年2月9日。已检索2月10日，2010。
 “新视野号，不完全是木星，首次发现冥王星”。pluto.jhuapl.edu - NASA新视野号任务现场。约翰霍普金斯大学应用物理实验室。2006年11月28日存档自原来的2014年11月13日。已检索11月29日，2011年。
 Chang，Kenneth (2016年10月28日)。"没有来自冥王星的更多数据"。纽约时报。已存档从原始的2019年3月29日。已检索10月28日，2016。
 冥王星探索完成: 新视野将2015年飞越数据的最后一部分返回地球。约翰霍普金斯大学应用研究实验室。2016年10月27日.已存档从原来的2016年10月28日。已检索10月28日，2016。
 布朗，德韦恩; 巴克利，迈克尔; 斯托托夫，玛丽亚 (2015年1月15日)。“发布15-011 - NASA的新地平线航天器开始冥王星遭遇的第一阶段”。NASA。已存档从原来的2020年4月7日。已检索1月15日，2015。
 “新视野”。pluto.jhuapl.edu。已存档从原来的2016年10月8日。已检索5月15日，2016。
 为什么一群科学家认为我们需要另一个冥王星任务。The Verge。已存档从原来的2018年7月8日。已检索7月14日，2018。
 为什么NASA应该再次访问冥王星。麻省理工学院技术评论。已存档从最初的2022年1月18日。已检索1月18日，2022年。
 “新视频模拟冥王星和Charon flyby; 提出了返回冥王星的任务”。2021年8月。存档自原来的2021年9月4日。已检索9月4日，2021。
 墙，迈克 (2017年5月8日)。“回到冥王星？科学家推动轨道飞行器任务“。Space.com。已存档从原来的2018年7月14日。已检索7月14日，2018。
 大厅，洛拉 (2017年4月5日)。“具有融合功能的冥王星轨道飞行器和着陆器”。NASA。已存档从原来的2017年4月21日。已检索7月14日，2018。
 具有融合功能的冥王星轨道飞行器和着陆器-第一阶段最终报告 已存档2019年4月29日，在回程机。(PDF) 斯蒂芬妮·托马斯，普林斯顿卫星系统公司。2017。
 Nadia Drake (2016年7月14日)。“我们在访问冥王星一年后学到的5件令人惊奇的事情”。国家地理。存档自原来的2021年3月7日。已检索8月19日，2021。
 “哈勃首次揭示冥王星的表面”。HubbleSite.org。太空望远镜科学研究所。1996年3月7日。已存档从原始的2021年8月19日。已检索10月18日，2021。
 "冥王星表面的地图"。HubbleSite.org。太空望远镜科学研究所。1996年3月7日。已存档从原始的2021年8月19日。已检索10月18日，2021。
 A.S.Ganesh (2021年3月7日)。“前所未有地看到冥王星”。印度教。已存档从原始的2021年8月19日。已检索8月19日，2021。
 Rothery，David A (2015年10月)。“来自新视野的冥王星和Charon”。天文学和地球物理学。56(5):5.19-5.22。doi:10.1093/astrogeo/atv168。
 劳尔，托德·R；约翰·R·斯宾塞；丹吉·伯特兰; 罗斯·A·拜尔；鲁尼翁，柯比·D；怀特，奥利弗·L；年轻，莱斯利A；金伯利·恩科; 威廉·B·麦金农；杰弗里·M·摩尔；凯瑟琳·B·奥尔金；斯特恩，S.艾伦; 韦弗，哈罗德·A。(2021年10月20日)。《冥王星的黑暗面》。行星科学杂志。2(214): 214。arXiv:2110.11976。Bibcode:2021PSJ.....2..214L。doi:10.3847/PSJ/ac2743。S2CID239047659。 
\end{enumerate}